\section{Downstream consequences: placement, roadmaps, and profiles}
\label{sec:consequence-pathways}

\begin{fixedbox}
\textbf{Downstream consequence-facing only.} From this point onward, the manuscript is packaging the already-closed theorem chain as downstream consequence material. Section~\ref{sec:consequence-pathways}, Section~\ref{sec:formalization-audit}, and the appendices record consequence placement, audit posture, source-role discipline, and atlas typing. None of this material reopens Sections~1--9 or adds a new foundational burden.
\end{fixedbox}

\subsection{Why application placement comes only after descriptive closure}

By the end of Section~\ref{sec:capstone-synthesis} the foundational theorem arc is already complete. The kernel floor is fixed, state-transition form and irreversible commitment are forced, the AMetric boundary is non-parameterized, the standing quotient and unique interior are fixed, mathematics and physics are closed under the same non-generative invariant, same-domain incompleteness objections are exhausted, and the manuscript has already reached the two-spectrum synthesis and the descriptive-closure theorem. What remains is therefore not a new foundational burden but a placement burden: how downstream manuscripts should appear without silently becoming proof sources for what is already closed.

\begin{importnote}[Downstream sources used here]
The application manuscripts are used here only in the exact roles their own front matter, theorem maps, and denial maps assign to them. \nsdoc{} is used for official-problem alignment, claim-class discipline, theorem-map hygiene, same-object conventions, one-route dependency audit, and theorem-local denial targets. \ymmaindoc{} and its three companions are used for endpoint hygiene, proof-home crosschecks, the theorem-level object of the local sharp-local net, the qualitative scope of faithful-Wilson universality, and the separation between local-net endpoints and nonlocal sector claims. \measurementdoc{} is used for the fixed-domain quotient theorem, the standing conditional, the confirmation criterion, the Everettian non-selection theorem, and the selector-exhaustion taxonomy. \structuradoc{} is used for witness-normal-form and transport-closure bookkeeping in the Standard Model case. None of these sources reopens Sections~1--9.
\end{importnote}

\subsection{Claim-type discipline}

\begin{definition}[Foundational closure theorem]
A \emph{foundational closure theorem} fixes same-scope admissibility, standing, non-generativity, quotient structure, or closure of admissible description at the level of the governing theorem architecture. Sections~1--9 belong to this class.
\end{definition}

\begin{definition}[Downstream internal problem solution]
A \emph{downstream internal problem solution} begins only after a closure theorem is already fixed and then uses that theorem to eliminate, transport, or classify the remaining licit endpoint statuses on a specific carrier. It may be theorem-heavy, but it does not function as the proof source of the foundational closure theorem from which it starts.
\end{definition}

\begin{definition}[Theorem roadmap and consequence profile]
A \emph{downstream consequence roadmap} is a compact main-body presentation of a downstream manuscript that records its exact claim type, theorem-level endpoint, live dependency spine, denial map, and non-claims. A \emph{downstream consequence profile} is a shorter placement statement used when the application clarifies the force of the closure theorem without needing its own extended theorem map inside the capstone. In this section, roadmap and profile are therefore both downstream consequence formats.
\end{definition}

\begin{theorem}[No downstream backflow]
\label{thm:no-downstream-backflow}
Let $M$ be a downstream application manuscript introduced after the foundational arc of Sections~1--9 is already complete. Under the declared source set, $M$ cannot function as a proof source for Sections~1--9 unless the theorem arc is explicitly rebuilt with $M$ promoted to a declared premise.
\end{theorem}

\begin{proof}
The foundational arc is already closed before any application manuscript enters. If a later application were allowed to backfill a premise of Sections~1--9 without explicit promotion, the manuscript would cease to have one strict dependency order. Therefore any such backflow must be blocked unless the theorem arc is rebuilt with that manuscript declared as a live premise.
\end{proof}

\begin{proposition}[Method-neutral reading of theorem roadmaps]
\label{prop:method-neutral-roadmaps}
A downstream manuscript may appear in the capstone as a theorem roadmap even when its proof class is not constructive, analytic, or estimate-driven in the conventional style of the target field. What matters for roadmap status is that the source manuscript itself fixes the official theorem object, states its exact claim class, identifies its live dependency spine, and localizes its own non-claims and denial map.
\end{proposition}

\begin{proof}
The capstone does not classify downstream manuscripts by conformity to a preferred proof style. It classifies them by theorem-role discipline. If a source manuscript states the official theorem object, fixes the method class actually claimed, exposes the live dependency spine, and marks what the manuscript does and does not prove, then the capstone can preserve that structure without importing the manuscript as a hidden premise. A theorem roadmap is therefore method-neutral in this exact sense: it records a theorem object and its audit architecture, not a stylistic preference for how such a theorem must be proved.
\end{proof}

\subsection{Navier--Stokes as a theorem roadmap}

\begin{proposition}[Why Navier--Stokes belongs here as a theorem roadmap]
Under the declared source set, the periodic Navier--Stokes manuscript should appear in the capstone as a theorem roadmap rather than as a proof-home appendix and rather than as a hidden foundational source. This remains true even though its proof class is explicitly a closure-induced impossibility theorem rather than a constructive continuation argument.
\end{proposition}

\begin{proof}
The manuscript \nsdoc{} is already claim-facing. Its opening states that the endpoint is Fefferman's periodic statement~(B), that the proof class is a \emph{closure-induced impossibility theorem} rather than a constructive continuation argument, that the theorem map is explicit, and that the official result is reached through one live route rather than through ambient corpus pressure. It also localizes the standard objections by theorem number and makes explicit what the result does \emph{not} claim: it does not assert that admissibility manufactures the missing analytic bridge inequality or a new continuation estimate; it asserts only that, once the official periodic classical scope and same-scope continuation question are fixed, any purported finite-horizon endpoint has no licit same-scope completion. Those are theorem-roadmap features rather than a mere application note, and they should be preserved in exactly that form here. In particular, the capstone should not rephrase the result as though its legitimacy depended on looking like a standard PDE continuation proof; the source manuscript itself already fixes the official theorem object and method-neutral claim class.\cite{MaleyCarrierNS}
\end{proof}

\begin{reviewbox}
\textbf{Imported hardening from the Navier--Stokes source.} The capstone should preserve six features from \nsdoc{} whenever it mentions the result: (i) the endpoint is the official periodic classical theorem, (ii) the proof class is closure-induced impossibility rather than constructive PDE continuation, (iii) the realized maximal classical solution is treated as one same-scope determination rather than as a sequence of independently reclassified time slices, (iv) no hidden endgame premises are allowed onto the live route, (v) the theorem-local denial map is part of the argument rather than a rhetorical add-on, and (vi) the same-scope status trichotomy is closed before the official theorem is compressed.\cite{MaleyCarrierNS}
\end{reviewbox}

\begin{proposition}[Method-neutral placement of the periodic Navier--Stokes manuscript]
The periodic Navier--Stokes manuscript retains theorem-roadmap status inside the capstone precisely because its source-level hardening already fixes three things theorem-locally: the official theorem object, the exact method class, and the one-route dependency audit. The capstone therefore preserves those features rather than translating the source into a conventional PDE-proof template.
\end{proposition}

\begin{proof}
The source manuscript states on its first pages that the endpoint is Fefferman's periodic statement~(B), that the proof class is a closure-induced impossibility theorem, and that the live route to the official theorem is unique and theorem-local. It also states what the manuscript does not claim: no new continuation estimate, no manufactured analytic bridge inequality, and no hidden endgame premise off the live route. Those declarations are not rhetorical. They are part of the theorem-local hardening. Preserving them is therefore the faithful capstone treatment.\cite{MaleyCarrierNS}
\end{proof}

\begin{theorem}[Closure--classical equivalence for the official periodic endpoint]
\label{thm:ns-closure-classical-equivalence}
Fix the official periodic scope: smooth divergence-free data on $\mathbb{T}^3$, fixed viscosity $\nu>0$, zero forcing, the fixed mean-zero pressure gauge, the lawful-redescription class generated by the standard periodic carrier skins, and the same-scope continuation question fixed by the source manuscript. Within that fixed scope, the following are equivalent.
\begin{enumerate}
  \item The official periodic classical endpoint holds: for every official datum $u_0\in C^\infty(\mathbb{T}^3;\mathbb{R}^3)$ with $\nabla\!\cdot u_0=0$, there exists a unique smooth periodic Navier--Stokes solution pair $(u,p)$ on $\mathbb{T}^3\times [0,\infty)$ with the fixed mean-zero pressure gauge.
  \item No licit same-scope finite-horizon endpoint status survives for a purported periodic breakdown. Equivalently: on the fixed periodic carrier, every purported finite-horizon endpoint proposal is either (i) represented by a standing-positive admissible in-scope witness, (ii) terminal standing failure, or (iii) an illicit same-scope gate / scope move; only clause~(i) is a licit same-scope endpoint status, and that remaining witness branch is then exhausted and eliminated.
\end{enumerate}
In particular, for the official periodic problem, the closure-induced impossibility result is not an adjacent interpretive claim but a theorem-local restatement of the official classical endpoint.
\end{theorem}

\begin{proof}
$(2)\Rightarrow(1)$ is exactly the manuscript-local live route compressed in the source text:
\[
19\mathrm{C}\text{--}19\mathrm{F}\ \to\ 20\mathrm{B}\ \to\ \mathrm{ST}4/\mathrm{ST}5\ \to\ 20\mathrm{C}\ \to\ 21.x.
\]
On that route, the fixed periodic problem first admits the unique admissible interior; the standard periodic carrier is identified with that interior; official data enter that carrier; the realized local class is shown to be exactly the full official local classical class; finite-horizon failure is already carrier-visible; the same-scope realized-lineage status trichotomy is proved and collapsed locally; and the remaining standing-positive witness branch is exhausted and eliminated. The source manuscript then states this official conclusion first as the transfer endpoint corollary and immediately afterward as the official classical restatement. $(1)\Rightarrow(2)$ is the theorem-local negative form of the same official endpoint: if the official periodic classical statement holds for the full official datum class, then no same-scope finite-horizon breakdown exists in that class, so no licit same-scope endpoint status survives for such a purported breakdown. Thus the two formulations are equivalent on the fixed official periodic scope.\cite{MaleyCarrierNS}
\end{proof}

\begin{theorem}[Exhaustiveness of the same-scope realized-lineage endpoint trichotomy]
\label{thm:ns-trichotomy-exhaustive}
Fix an admitted carrier state $x$ for the official periodic problem and let $R_{\mathrm{NS}}(x)=(u_x,p_x)$ be its realized maximal classical lineage on $[0,T^*(x))$. Assume $T^*(x)<\infty$. Then every purported same-scope finite-horizon endpoint proposal for this realized periodic lineage has exactly one of the following theorem-local statuses:
\begin{enumerate}
  \item[(i)] it is represented by a standing-positive admissible in-scope witness on the fixed torus carrier;
  \item[(ii)] it is terminal standing failure; or
  \item[(iii)] it is an illicit same-scope gate / scope move, namely a path-dependent reclassification, representative-only or metric-only gate, new selector / tolerance / discriminator, or undeclared scope change.
\end{enumerate}
No fourth same-scope endpoint status exists.
\end{theorem}

\begin{proof}
The realized maximal classical periodic object is treated as one same-scope determination with one fixed continuation question, not as a sequence of independently reclassified time slices. On that fixed same-scope lineage, admissibility and standing are global rather than time-indexed, and every standing-relevant same-scope discriminator factors through the future-behavior quotient. Therefore an endpoint proposal cannot survive as a fresh stage-indexed status, a path-dependent status, or a representative-only / skin-level status.

If the endpoint is presented through the carrier-side standing-relevant horizon witness, it lies in clause~(i). If instead the proposal says that the same realized lineage can be completed only by leaving admissible representability while still claiming the same torus carrier and lineage identity, that is exactly terminal standing failure, giving clause~(ii). If it tries to avoid both routes while keeping the endpoint alive, it must smuggle in a new same-scope discriminator, representative-only or metric-only gate authority, selector, tolerance, or undeclared scope change, giving clause~(iii). These cases are exhaustive because the same-object convention, quotient / history-factorization packet, no-admissibility-dynamics theorem, and ATS non-authority result leave no additional same-scope way to distinguish or preserve an endpoint claim. Hence no fourth same-scope endpoint status exists.\cite{MaleyCarrierNS}
\end{proof}

\subsection{Yang--Mills as a theorem roadmap}

\begin{proposition}[Why Yang--Mills belongs here as a theorem roadmap]
Under the declared source set, the Yang--Mills cluster should appear in the capstone as a theorem roadmap rather than as a proof-home appendix and rather than as a hidden foundational source.
\end{proposition}

\begin{proof}
The Yang--Mills endpoint manuscript \ymmaindoc{} already states the theorem-level object, the exact claim boundaries, the theorem-map reduction, the seam roles of the three companions, and the local-net endpoint. It fixes from the outset that the theorem-level object is the bounded-region local gauge-invariant sharp-local net generated by flowed curvature composites together with its OS/Wightman vacuum representation and vacuum-sector mass gap; that faithful-Wilson universality is qualitative only within the admissible gradient-flow renormalization class and after coherent normalization; that the explicit quantitative Route~1 ledger remains $\rho$-indexed; and that extended-support line/surface sectors remain outside the theorem-level claim. The three companions then preserve fixed role splits rather than adding theorem-level ambiguity: Companion~I isolates the ultraviolet gate and the live Route~1 mass-gap chain, Companion~II isolates the Lane~A sharp-local construction and the finite-cap / inductive-union seam discipline, and Companion~III isolates reconstruction, non-triviality, faithful-Wilson universality, and the exact local-net endpoint. That is theorem-roadmap architecture, and the capstone should preserve it in exactly that form rather than compressing it into a generic existence narrative.\cite{MaleyYangMillsNet,MaleyYMCompanionI,MaleyYMCompanionII,MaleyYMCompanionIII}
\end{proof}

\begin{reviewbox}
\textbf{Imported hardening from the Yang--Mills sources.} The capstone should preserve five items from the Yang--Mills cluster: (i) the theorem-level object is local and gauge-invariant, (ii) faithful-Wilson universality is qualitative and normalization-bound rather than an unqualified quantitative uniformity claim, (iii) the Route~1 ledger is explicitly $\rho$-indexed, (iv) the companions carry fixed proof-home role splits, and (v) nonlocal sectors are excluded from the theorem-level endpoint rather than silently absorbed into it.\cite{MaleyYangMillsNet,MaleyYMCompanionI,MaleyYMCompanionII,MaleyYMCompanionIII}
\end{reviewbox}

\begin{proposition}[Endpoint-hygiene placement of the Yang--Mills roadmap]
The Yang--Mills materials retain theorem-roadmap status inside the capstone because the source cluster already fixes the theorem-level object, the exact endpoint boundaries, the role of qualitative faithful-Wilson universality, the $\rho$-indexed status of the explicit quantitative ledger, and the exclusion of nonlocal sectors from the theorem claim.
\end{proposition}

\begin{proof}
The endpoint manuscript and the three companions divide these burdens explicitly. The main endpoint text fixes the theorem-level local sharp-local net and the vacuum-sector mass gap, while the companions freeze the ultraviolet gate, Route~1 chain, Lane~A seam discipline, reconstruction, non-triviality, and faithful-Wilson universality. Because these role splits are already source-level and theorem-facing, the capstone should preserve them as a roadmap rather than repackage them as a single undifferentiated application note.\cite{MaleyYangMillsNet,MaleyYMCompanionI,MaleyYMCompanionII,MaleyYMCompanionIII}
\end{proof}

\subsection{Measurement as a compact consequence profile}

\begin{proposition}[Why the measurement manuscript belongs here as a consequence profile]
Under the declared source set, the measurement-reference manuscript should appear in the capstone as a compact consequence profile rather than as a theorem roadmap and not as a foundational proof source.
\end{proposition}

\begin{proof}
The manuscript \measurementdoc{} already carries a complete fixed-domain theorem ladder: admissible reference assignments, the measurement-reference quotient theorem, global equivalence collapse, the Everettian non-selection theorem, the fixation-commutation criterion for theory-level confirmation, diachronic inheritance for memory, anticipation, deliberation, and confirmation, and selector exhaustion. It also states the exact target discourse and its denial-condition matrix. Those hardenings are real, but the application is narrower than the foundational arc and is explicitly conditional on the downstream closure package already fixed earlier. Its role in the capstone is therefore to show what the closure theorem does to a physically non-selective measurement theory, not to introduce a second theorem spine in full. The correct placement is a compact consequence profile stating the claim type, the standing conditional, and the exact denial targets without reproducing the entire measurement manuscript.\cite{MaleyMeasurement}
\end{proof}

\begin{reviewbox}
\textbf{Imported hardening from the measurement source.} The capstone preserves six source-level features whenever it cites \measurementdoc{}: (i) the standing conditional that fixes the target explanatory practice, (ii) the measurement-reference quotient theorem, (iii) global equivalence collapse on the same scope, (iv) the Everettian non-selection theorem for the bare physical package, (v) the fixation-commutation criterion for theory-level confirmation, and (vi) the selector-exhaustion taxonomy together with the denial-condition matrix.\cite{MaleyMeasurement}
\end{reviewbox}

\subsection{Standard Model as a compact consequence profile}

\begin{proposition}[Why the Standard Model manuscript belongs here as a consequence profile]
Under the declared source set, the Standard Model bookkeeping manuscript should appear in the capstone as a compact consequence profile rather than as a theorem roadmap.
\end{proposition}

\begin{proof}
The role of \structuradoc{} in the capstone is explanatory rather than proof-bearing. It gives a witness-normal-form and transport-closure ledger for a fixed renormalizable chiral gauge load, but it does not reopen the theorem class of Sections~1--9. Its contribution is to show, in a concrete physical domain, how standing-preserving redescription and transport closure collapse apparent parameter freedoms and certify which inputs remain irreducible. That is exactly the work of a compact consequence profile.\cite{MaleyStructura}
\end{proof}

\subsection{Exact placement theorem}

\begin{theorem}[Main-body placement theorem]
\label{thm:application-placement}
Under the declared source set, the main body of the capstone should present downstream applications in the following exact form:
\begin{enumerate}
  \item the periodic Navier--Stokes and Yang--Mills manuscripts appear as theorem roadmaps;
  \item the measurement-reference and Standard Model manuscripts appear as compact consequence profiles;
  \item a Zenodo-only appendix atlas may extend the downstream layer by adding further theorem roadmaps or compact consequence profiles, provided each appendix is consequence-facing only, explicitly conditional on Sections~1--9, and removable without altering the foundational arc; and
  \item none of these downstream manuscripts is a live dependency of the foundational closure arc already completed in Sections~1--9.
\end{enumerate}
\end{theorem}

\begin{proof}
The first clause follows from the theorem-map, claim-type, denial-map, and endpoint-hygiene hardening already present in \nsdoc{} and in the Yang--Mills cluster. The second follows from the narrower, downstream, and explicitly conditional role of \measurementdoc{} and \structuradoc{}. The third clause is the no-backflow theorem applied to a controlled appendix atlas: once additional materials are typed as theorem roadmaps or compact consequence profiles and are declared removable without changing Sections~1--9, they extend explanatory coverage without becoming proof sources. The fourth clause is Theorem~\ref{thm:no-downstream-backflow} itself.
\end{proof}

\subsection{Zenodo appendix-atlas extension}

\begin{proposition}[Zenodo appendix-atlas rule]
\label{prop:zenodo-appendix-atlas-rule}
A Zenodo consequence atlas may be appended after the audit architecture if and only if every appendix satisfies all of the following conditions:
\begin{enumerate}
  \item it is typed explicitly as a theorem roadmap or a compact consequence profile;
  \item it declares its fixed-domain object and the exact sections of the foundational arc on which it depends;
  \item it states an exact no-backflow clause; and
  \item deleting the appendix leaves Sections~1--9 unchanged.
\end{enumerate}
\end{proposition}

\begin{proof}
Each condition is forced by the claim-type discipline already fixed in this section. Without explicit typing, a downstream appendix drifts toward hidden-proof-source status. Without a fixed-domain declaration and dependency note, scope drift is no longer controlled. Without a no-backflow clause, the appendix can be misread as silently rebuilding an earlier theorem burden. And if deleting the appendix changes Sections~1--9, then the appendix was never merely consequence-facing. Therefore the four conditions are necessary and jointly sufficient for a Zenodo-only appendix atlas compatible with the capstone architecture.
\end{proof}

\begin{fixedbox}
\textbf{Exact wording for the 42 version.} The capstone proves the foundational closure theorem chain in Sections~1--9. Navier--Stokes and Yang--Mills appear in the main body as theorem roadmaps, and measurement and Standard Model materials appear there as compact consequence profiles. The Zenodo-only appendix atlas extends the downstream layer across further controlled theorem roadmaps and consequence profiles, each consequence-facing only and each removable without altering the foundational arc. None of the appendix manuscripts functions as a premise of Sections~1--9.
\end{fixedbox}
